\chapter*{Prefacio}
\addcontentsline{toc}{chapter}{Prefacio}

\begin{flushleft}\small
\textbf{Lo que vas a leer.}\\[0.3em]
Este libro es la guía pedagógica de los algoritmos, modelos y trucos
de sistemas que hacen que Graph Hunter ---el motor de \emph{threat
hunting} basado en grafos temporales que desarrollamos en
Base4Security--- corra sub-10\,ms sobre grafos de 100K aristas y
sobreviva sesiones de 28\,GB de EVTX en una notebook. La meta no es
documentar el código (de eso se ocupan los libros \emph{Arquitectura}
y \emph{Especificación}), ni enseñar a operarlo (\emph{Manual}), ni
defenderlo académicamente (\emph{Paper}). La meta es que un lector
con \emph{background} de CS de primer o segundo año pueda entender,
desde cero, \emph{por qué} cada decisión es la que es.
\end{flushleft}

\section*{¿A quién está dirigido?}

A tres audiencias bien delimitadas:

\begin{itemize}
    \item \textbf{Estudiantes de Ciencias de la Computación} que ya
    saben qué es un grafo dirigido, big-O, recursión, \emph{hash
    tables}, y han programado en algún lenguaje de sistemas (C, C++,
    Rust o Go). Para vos este libro es un \emph{textbook} de un
    sistema real: no es un toy-example didáctico, es la mecánica
    interna de un motor que sostiene casos de fraude reales.
    \item \textbf{Ingenieros de \emph{platforms} de seguridad}
    (\emph{detection engineers}, \emph{SIEM} engineers, \emph{XDR}
    devs) que vienen de full-text search y quieren entender qué cambia
    cuando el modelo de datos pasa de tabla a grafo temporal.
    \item \textbf{Devs nuevos al equipo de Graph Hunter} que necesitan
    ramp-up rápido sin tener que leer 50.000 líneas de Rust antes de
    proponer una optimización. Este libro es la versión larga de la
    onboarding chat.
\end{itemize}

\section*{Prerequisitos}

Asumimos que:

\begin{itemize}
    \item Sabés qué es un \emph{grafo dirigido} y un \emph{árbol},
    distinguís \emph{path} de \emph{walk}, y entendés \emph{depth-first
    search} y \emph{breadth-first search} a nivel pseudocódigo.
    \item Manejás notación big-O y entendés que ``\(O(m \cdot
    \deg^k)\)'' significa algo concreto sobre cómo escala un algoritmo.
    \item Conocés \emph{hash tables} (\(O(1)\) amortizado, función de
    hash, colisiones) y \emph{arrays}.
    \item Programaste alguna vez en un lenguaje de sistemas y sabés
    qué es un puntero, qué pasa cuando se desreferencia, y por qué
    asignar memoria no es ``gratis''.
\end{itemize}

\noindent\textbf{No} asumimos:

\begin{itemize}
    \item Conocimiento previo de bases de datos relacionales o álgebra
    relacional. Cuando aparezca un \emph{join}, lo definimos
    (capítulo~\ref{cap:joins-y-grafos}).
    \item Familiaridad con \emph{worst-case-optimal joins}, AGM bound,
    Yannakakis, GYO, hipergrafos. Construimos cada uno desde cero.
    \item Conocimiento de \emph{memory-mapped files}, \emph{cache
    lines}, \emph{branch prediction}. El capítulo~\ref{cap:sistemas}
    los cubre antes de que los necesites.
    \item Familiaridad con Rust. Los \emph{listings} de código son
    auto-explicativos para alguien acostumbrado a un \emph{C-family
    language}; cuando una construcción de Rust ($\langle$\,\texttt{Arc},
    \texttt{RwLock}, \texttt{repr(C)}\,$\rangle$) sea relevante para el
    argumento, la explicamos en el momento.
\end{itemize}

\section*{Cómo está organizado el libro}

\begin{description}
    \item[Parte~I — Fundamentos.] Cinco capítulos que cubren todo lo
    que necesitás saber \emph{para poder leer el resto del libro}: el
    paradigma de threat hunting, la formalización de grafos
    temporales, el problema de subgraph matching, los joins
    relacionales (con derivación didáctica del AGM bound), y los
    fundamentos de sistemas que justifican las decisiones de
    \emph{layout} y memoria del motor.
    \item[Partes II--III — Arquitectura y modelo de datos.] Tres
    capítulos de arquitectura (capas, pipeline) y cuatro del modelo
    de datos: \texttt{Entity}, \texttt{Relation}, el \emph{string
    interner}, \texttt{CompactRelation} de 32 bytes, y la
    migración del \texttt{HuntResult} a \texttt{Vec<Arc<str>>}.
    \item[Parte~IV — Pipeline de ingesta.] De \emph{log} a grafo
    temporal: parsers, \emph{streaming mirror}, \emph{lazy build} de
    índices, NLF incremental.
    \item[Parte~V — El \emph{matcher}, paso a paso.] El núcleo
    algorítmico del libro. Diez capítulos que construyen el zoológico
    de operadores (DFS legacy con \emph{prunes}, LFTJ-K\(_4\),
    Yannakakis chain, Yannakakis tree-walk) desde el DFS recursivo
    más básico hasta los \emph{worst-case-optimal joins}.
    \item[Parte~VI — Índices auxiliares.] NLF y \emph{k-hop
    reachability} en profundidad: \emph{layouts} \emph{dense}+\emph{sparse},
    \emph{mmap-spill}, \emph{lazy build}, extensión temporal.
    \item[Parte~VII — La DSL.] Diseño de una \emph{domain-specific
    language} chica: forma \emph{chain} legacy, forma \emph{graph} con
    el separador \texttt{;}, parser, AST, dispatch al \emph{matcher}.
    \item[Parte~VIII — Diseño futuro (L3).] Las dos piezas pendientes
    de la \emph{roadmap} L3: \emph{delta} \emph{khop\_reach} usando
    propagación inversa estilo Kankanamge SIGMOD 2017, y
    \emph{streaming} como \emph{canonical store}.
    \item[Parte~IX — Plug-and-play.] La tesis transversal: cada
    optimización que entrega valor lo hace sin que el usuario tenga
    que tocar \emph{config}. Cómo se materializa esa promesa en el
    \emph{frontend} y la \emph{telemetría}.
    \item[Apéndices.] Notación, glosario, soluciones a ejercicios
    selectos, bibliografía.
\end{description}

\section*{Cómo leer este libro}

\begin{itemize}
    \item \textbf{Lectura lineal.} Si no conocés el motor, leélo de
    principio a fin. Cada capítulo asume los anteriores; los
    \emph{cross-refs} son explícitos cuando un argumento depende de
    algo de otro capítulo.
    \item \textbf{Lectura por reference.} Si ya sabés de qué se trata
    y querés entender un componente concreto, andá directo al
    capítulo correspondiente. Cada uno arranca con una caja
    \emph{Intuición} que resume el contenido en un párrafo, y termina
    con una sección \emph{Ejercicios}; podés saltearte el cuerpo y
    volver a él si la intuición no alcanza.
    \item \textbf{Lectura para enseñar.} Cada capítulo está pensado
    para ser una clase de 50--90~minutos. Los ejercicios marcados
    con \starred{} tienen solución en el Apéndice~C; los demás son
    abiertos para que el lector se trabe un poco antes de pedir
    ayuda.
\end{itemize}

\section*{Convenciones tipográficas}

\begin{itemize}
    \item Las cajas \textbf{Intuición} (verde claro) preceden a la
    formalización. Si entendés la intuición, podés saltearte la prueba.
    \item Las cajas \textbf{Ejemplo} (borde verde claro, fondo blanco)
    contienen ejemplos trabajados extensos paso a paso.
    \item \texttt{Texto monoespaciado} es código o nombre de tipo /
    función / archivo del repo.
    \item \emph{Texto en cursiva} marca términos técnicos en su primera
    aparición o préstamos lingüísticos del inglés (\emph{path},
    \emph{join}, \emph{spill}).
    \item \(\mathcal{V}\), \(\mathcal{E}\), \(\mathcal{T}\),
    \(\mathcal{L}\) son los conjuntos de vértices, aristas, tiempo y
    etiquetas respectivamente. Convenciones completas en
    Apéndice~\ref{ap:notacion}.
    \item \texttt{\filepath{core/graph-engine/src/lftj.rs}} (con
    fuente monoespaciada) cita un archivo del repo. Las cajas
    \textbf{En el código} apuntan del concepto al fichero real cuando
    eso ayuda al lector a auditar la implementación.
    \item Los \emph{algorithm blocks} usan pseudocódigo estilo CLRS:
    \texttt{for}, \texttt{if}, \texttt{while}, indentación por bloque,
    asignación con \(\leftarrow\). Convenciones completas en la
    sección \emph{Convenciones} (a continuación) y en el
    Apéndice~\ref{ap:notacion}.
\end{itemize}

\section*{Relación con los otros libros}

GraphHunter tiene seis libros. Cada uno cubre una porción distinta:

\begin{center}
\renewcommand{\arraystretch}{1.2}
\begin{tabular}{p{2.4cm}p{10.5cm}}
\toprule
\textbf{Libro} & \textbf{Cubre} \\
\midrule
\emph{Algoritmos} & (este libro) — explicación pedagógica desde lo básico, con ejemplos y ejercicios. \\
\emph{Diseño} & resumen ejecutivo del diseño en \(\sim\,25\) páginas, asumiendo que el lector ya conoce el dominio. \\
\emph{Arquitectura} & estructura de \emph{crates}, \emph{layering}, ADRs, decisiones de \emph{repo organization}. \\
\emph{Especificación} & \emph{wire contracts} (HTTP API, MCP, FFI, formatos de \emph{log}, \emph{env vars}). \\
\emph{Manual} & guía operacional para el analista que usa el motor. \\
\emph{Paper} & justificación académica con \emph{benchmarks} y comparaciones contra el estado del arte. \\
\bottomrule
\end{tabular}
\end{center}

Si una decisión te resulta opaca acá, casi seguro la otra cara
(\emph{wire contract}, \emph{benchmark}, ADR) está en uno de los otros
cinco libros, y el \emph{cross-ref} explícito te lleva ahí.

\vspace{1em}
\begin{flushright}
\textit{Lucas Sotomayor}\\
\textit{Base4Security}\\
\textit{Mayo 2026}
\end{flushright}

\cleardoublepage
